\section{Pengaruh Kedalaman Sirkuit (\textit{Depth}) Terhadap Akurasi Solusi}
\label{sec:4_8_depth_analysis}

Eksperimen terhadap variasi kedalaman sirkuit (\textit{depth}) dilakukan untuk mengevaluasi pengaruh kompleksitas ansätze terhadap kemampuan algoritma dalam menemukan energi minimum Hamiltonian. Secara teoretis, peningkatan \textit{depth} akan menambah jumlah parameter bebas dan gerbang keterikatan, yang secara langsung meningkatkan \textit{expressibility} sirkuit untuk memetakan fungsi gelombang yang lebih presisi. Berdasarkan grafik pada \texttt{grafik\_depth\_per\_window\_N2.png}, terlihat bahwa peningkatan kedalaman dari 1 hingga 6 lapis cenderung memberikan hasil energi yang lebih rendah secara konsisten. Hal ini menunjukkan bahwa sirkuit yang lebih dalam memiliki fleksibilitas lebih besar untuk menyelaraskan status kuantum dengan konfigurasi \textit{ground state} yang diinginkan.

Namun, pengamatan mendalam mengungkap adanya titik jenuh (\textit{saturation point}) di mana penambahan kedalaman sirkuit tidak lagi memberikan penurunan energi yang signifikan. Pada sistem dengan jumlah qubit kecil ($N=2$), penggunaan \textit{depth} yang terlalu tinggi justru dapat menjadi kontraproduktif karena meningkatkan beban komputasi dan risiko \textit{over-parameterization}. Peningkatan jumlah parameter yang tidak perlu seringkali memicu fenomena \textit{barren plateaus} yang lebih parah, di mana lanskap energi menjadi sangat datar sehingga \textit{optimizer} SPSA memerlukan waktu iterasi yang jauh lebih lama untuk menemukan arah gradien yang benar. Oleh karena itu, pemilihan kedalaman sirkuit yang optimal harus menyeimbangkan antara kebutuhan akan akurasi solusi dengan efisiensi penggunaan sumber daya komputasi.

Analisis terhadap korelasi antara \textit{depth} dan dinamika entropi juga menunjukkan bahwa sirkuit yang lebih dalam cenderung menghasilkan tingkat \textit{entanglement} residu yang lebih tinggi. Seiring bertambahnya jumlah gerbang keterikatan, sistem menjadi semakin sulit untuk diruntuhkan (\textit{collapse}) kembali menjadi status klasik murni di akhir iterasi. Hal ini menjelaskan mengapa pada \textit{depth} yang tinggi, entropi seringkali gagal menyentuh angka nol meskipun energi yang dicapai sudah cukup rendah. Temuan ini mengindikasikan bahwa untuk permasalahan portofolio skala kecil, struktur sirkuit yang moderat (sekitar \textit{depth} 2 atau 3) seringkali lebih unggul dalam memberikan solusi yang tidak hanya berenergi rendah, tetapi juga memiliki tingkat kepercayaan probabilitas yang lebih tinggi dan deterministik.
